\chapter{Linear Model under General Variance}

\section{Model}

Consider the linear model
\[
y = X\beta + e,
\]

In the classical linear model, we assumed
\[
V(e)=\sigma^2 I_T.
\]

\paragraph{Intuition}
This assumption means that all error terms have:
\begin{itemize}
    \item the same variance (homoskedasticity),
    \item zero covariance across observations (no correlation).
\end{itemize}

In reality, this is often too restrictive. For example:
\begin{itemize}
    \item Some observations may be noisier than others (heteroskedasticity).
    \item Time series data often exhibit correlation across time (autocorrelation).
\end{itemize}

Now we relax this assumption and allow
\[
V(e)=\sigma^2 \Psi,
\]

where $\Psi$ is a symmetric positive definite matrix.

\paragraph{Interpretation}
The matrix $\Psi$ captures the full variance structure:
\begin{itemize}
    \item diagonal elements = variances of each observation,
    \item off-diagonal elements = covariances between observations.
\end{itemize}

Thus, $\Psi$ tells us how “reliable” each observation is and how they relate to each other.

---

\section{Transformation of the Model}

Since $\Psi$ is positive definite, $\Psi^{-1}$ exists. Let
\[
\Psi^{-1} = P'P.
\]

\paragraph{Key Idea}
We want to transform the model so that the new error term satisfies the classical assumptions.

Premultiply the model:
\[
Py = PX\beta + Pe.
\]

Define
\[
y^* = Py, \quad X^* = PX, \quad e^* = Pe.
\]

Then
\[
y^* = X^*\beta + e^*.
\]

Now compute the variance:
\[
V(e^*) = P V(e) P' = P(\sigma^2 \Psi)P' = \sigma^2 P\Psi P'.
\]

Since
\[
\Psi^{-1} = P'P,
\]
we obtain
\[
P\Psi P' = I_T.
\]

Therefore,
\[
V(e^*) = \sigma^2 I_T.
\]

\paragraph{Conclusion}
We have transformed the model into a standard classical linear model.

\paragraph{Deep Intuition}
The transformation matrix $P$ “whitens” the data:
\begin{itemize}
    \item it removes correlation,
    \item it rescales observations to equal variance.
\end{itemize}

Thus,
\[
\text{GLS = OLS applied to transformed data}
\]

---

\section{Generalized Least Squares}

The transformed model implies minimizing:
\[
S = e^{*'} e^*.
\]

Since $e^* = Pe$, we have:
\[
S = e'P'Pe = e'\Psi^{-1}e.
\]

Thus,
\[
S(\beta)
=
(y-X\beta)'\Psi^{-1}(y-X\beta).
\]

\paragraph{Interpretation}
This is called the \textbf{weighted least squares} criterion.

\paragraph{Key Insight}
Unlike OLS, GLS assigns different weights:
\begin{itemize}
    \item Observations with large variance → lower weight
    \item Observations with small variance → higher weight
\end{itemize}

\paragraph{Why $\Psi^{-1}$?}
Because the inverse variance determines optimal weighting:
\[
\text{weight} \propto \frac{1}{\text{variance}}
\]

---

First-order condition:
\[
X'\Psi^{-1}X\beta = X'\Psi^{-1}y.
\]

Thus,
\[
\hat{\beta}_{GLS}
=
(X'\Psi^{-1}X)^{-1}X'\Psi^{-1}y.
\]

---

\section{Properties of GLS}

\subsection{Unbiasedness}

\[
E(\hat{\beta}_{GLS})=\beta.
\]

\paragraph{Insight}
GLS remains unbiased because it is still a linear function of $y$ and $E(e)=0$.

---

\subsection{Variance}

\[
V(\hat{\beta}_{GLS})
=
\sigma^2 (X'\Psi^{-1}X)^{-1}.
\]

\paragraph{Key Point}
This variance explicitly accounts for heteroskedasticity and correlation.

---

\subsection{Efficiency}

GLS is BLUE.

\paragraph{Deep Intuition}
OLS ignores $\Psi$, while GLS uses it.

Thus:
\begin{itemize}
    \item OLS = correct mean, wrong weighting
    \item GLS = correct mean, correct weighting
\end{itemize}

Therefore,
\[
V(GLS) \leq V(OLS).
\]

---

\section{Comparison between OLS and GLS}

OLS:
\[
\hat{\beta}_{OLS}=(X'X)^{-1}X'y
\]

GLS:
\[
\hat{\beta}_{GLS}=(X'\Psi^{-1}X)^{-1}X'\Psi^{-1}y
\]

\paragraph{Core Difference}
OLS treats all observations equally, while GLS adjusts for reliability.

\paragraph{Geometric Insight}
OLS minimizes Euclidean distance:
\[
e'e
\]

GLS minimizes a weighted distance:
\[
e'\Psi^{-1}e
\]

---

\subsection{OLS under General Variance}

\[
E(\hat{\beta}_{OLS})=\beta
\]

\[
V(\hat{\beta}_{OLS})
=
\sigma^2 (X'X)^{-1}X'\Psi X(X'X)^{-1}.
\]

\paragraph{Important Insight}
OLS is still unbiased but inefficient.

\paragraph{Interpretation}
Ignoring $\Psi$ leads to suboptimal use of information.

---

\section{Estimation of $\sigma^2$}

\[
s^2_{GLS}
=
\frac{
(y-X\hat{\beta}_{GLS})'\Psi^{-1}(y-X\hat{\beta}_{GLS})
}{T-K}.
\]

\paragraph{Key Idea}
Residuals must also be weighted consistently with the model.

---

\section{Prediction}

\[
y_0 = X_0\beta + e_0.
\]

\[
\hat{y}_0
=
X_0\hat{\beta}_{GLS}
+
C'\Psi^{-1}(y-X\hat{\beta}_{GLS}).
\]

\paragraph{Intuition}
Prediction combines:
\begin{itemize}
    \item structural component: $X_0\beta$
    \item correction term: based on past errors
\end{itemize}

\paragraph{Special Case}
If $C=0$, then
\[
\hat{y}_0 = X_0\hat{\beta}_{GLS}.
\]

\paragraph{Meaning}
If errors are uncorrelated across samples, no correction is needed.

---

\section{Hypothesis Testing}

\[
H_0: R\beta = r
\]

Test statistic:
\[
F
=
\frac{(WSSE_R-WSSE_u)/J}{WSSE_u/(T-K)}.
\]

\paragraph{Interpretation}
Same structure as OLS F-test, but:
\begin{itemize}
    \item SSE replaced with weighted SSE
\end{itemize}

---

\section{Feasible GLS}

\paragraph{Problem}
GLS requires knowing $\Psi$, which is usually unknown.

\paragraph{Solution}
Estimate $\Psi$ from data.

\[
\hat{\beta}_{FGLS}
=
(X'\hat{\Psi}^{-1}X)^{-1}X'\hat{\Psi}^{-1}y.
\]

\paragraph{Key Insight}
FGLS approximates GLS using estimated variance structure.

---

\section{FGLS Procedure}

\begin{enumerate}
    \item Estimate OLS:
    \[
    \hat{\beta}_{OLS}
    \]
    \item Compute residuals:
    \[
    \hat{e}
    \]
    \item Estimate $\Psi$
    \item Apply GLS using $\hat{\Psi}$
\end{enumerate}

\paragraph{Big Picture}
\[
\text{OLS} \rightarrow \text{estimate structure} \rightarrow \text{GLS}
\]

---

\section{Asymptotic Properties of FGLS}

If
\[
\hat{\Psi} \xrightarrow{p} \Psi,
\]

then
\[
\hat{\beta}_{FGLS}
\approx
N\left(
\beta,
\sigma^2 (X'\Psi^{-1}X)^{-1}
\right).
\]

\paragraph{Important Insight}
FGLS behaves like GLS in large samples.

---

\section{Final Summary}

\begin{itemize}
    \item OLS is always unbiased
    \item OLS is inefficient if $\Psi \neq I$
    \item GLS is the optimal estimator
    \item FGLS is the practical implementation
\end{itemize}

\[
\text{OLS: simple but inefficient}
\]
\[
\text{GLS: optimal but requires }\Psi
\]
\[
\text{FGLS: feasible approximation}
\]